\section{From Field to Watershed: Mapping Networks and Scaling Response}

\subsection{Overview}

Two linked spatial problems anchor this frontier. First, the physical drainage
network---where tiles lie, at what depth and spacing, and how mains connect
fields to ditches---is largely unmapped, so uncertainty about the network
itself propagates into every downstream hydrologic and water-quality model.
Second, even where field processes are well understood, the field is not the
scale at which water-quality goals are set; translating instrumented-plot
knowledge to the watersheds where nutrient and sediment targets apply remains a
genuinely open methodological problem. The two are related: because networks
are unknown and because routing, connectivity, and in-stream processes are
weakly represented, effectiveness estimates for practices such as drainage
water management and phosphorus mitigation cannot yet be aggregated with
confidence from field to basin. The priorities that recur across the literature
are therefore concrete---non-invasive detection methods with usable
practitioner guidelines, shared drainage databases with standardised reporting,
intensively monitored small catchments as upscaling bridges, and mechanistic
treatment of the channel-network, connectivity, and depressional processes that
mediate the field-to-watershed link.

\subsection{Locating, mapping, and cataloguing drainage networks}

Because the drainage network is often unknown, detection and mapping form a
fast-moving frontier in which methods work under some conditions but lack
systematic guidelines. Actual tile locations and measured discharge would
sharpen model calibration \parencite{Amado2018}, and predictive mapping from
imagery needs validation against GPS surveys plus methods that avoid
misidentifying noise and that reconnect broken segments
\parencite{Naz2009,Schilling2019}. The most explicit data agenda is
institutional: report tile diameter, depth, spacing, and mains, and build a
shared drainage database jointly maintained by agencies, researchers,
landowners, and contractors, with repositories updated as records accrue
\parencite{Zhou2025,Valayamkunnath2020}.

\subsubsection{Remote sensing and machine learning}

Remote-sensing work seeks urgent development of non-invasive methods for aging
systems, including UAV flights combined with ground-penetrating radar
\parencite{Wienken2024}, and multispectral, thermal-infrared, and
machine-learning agendas that define the site conditions under which drain
lines and outlets can be mapped and that yield decision-tree guidelines for
practitioners \parencite{Allred2020,Allred2021,Allred2022}. \textcite{Carlsen2024}
recommend reconsidering the often-cited 2.54~cm minimum rainfall for imagery
acquisition and investigating deep learning, multitemporal acquisition, and
synthetic aperture radar while expanding public drainage-map datasets---echoed
by calls for transformer-based architectures and higher-resolution imagery
integrated into hydrologic frameworks \parencite{Wan2024} and for reduced
tracing uncertainty through multiple interpreters, ground truth, and richer
covariates \parencite{Redoloza2023,Harris2018,Yimer2023}. Optical and radar
covariate studies further seek more frequent imaging, repeated SAR with ground
measurements, and UAV cover-crop mapping
\parencite{Tlapáková2017,Meunier2001,Xu2018}.

\subsubsection{Geophysics}

Geophysical evaluations remain condition-dependent and call for data from more
sites and soils \parencite{Allred2005}, systematic testing of GPR antenna
configuration and orientation \parencite{Allred2013}, and prediction of GPR
suitability via electromagnetic simulation driven by EMI. Concurrent EMI/GPR
surveys and drone or magnetic surveys are recommended, with performance shown to
depend on soil, residue, tillage, and wetness
\parencite{Koganti2020,Koganti2021,Jouen2018,Kratt2020,Allred2018b}. The
open question is less whether any single method can detect drains than how to
specify, in advance, which method suits a given soil, cover, and moisture
state.

\subsection{Scaling from field to watershed}

The defining open problem is the disconnect between field-scale process
knowledge and watershed-scale response. Quantifying the relationships between
field processes and watershed response is repeatedly identified as critical for
water-quality strategy, motivating explicit field-to-watershed linkage and
upscaling frameworks \parencite{Williams2015,Williams2017,Wang2020b}. Small,
intensively monitored catchments are proposed as the empirical bridge for
upscaling and model validation \parencite{Kyllmar2023,Basu2010}, and
effectiveness estimates for drainage water management and phosphorus mitigation
should be scaled up from field to watershed
\parencite{Wang2021,Ross2016,Shedekar2021}.

The hydrologic effects of tiling remain inadequately characterised at watershed
scale and must be pursued through combined modelling and field data: attending
to seasonal streamflow recessions in tiled versus non-tiled watersheds
\parencite{Schilling2008,Schilling2012}, to events of all sizes across
landscapes \parencite{Sloan2017}, to routing between spatially variable
depressional stores and flow accumulation across river networks
\parencite{Sloan2016,Fonley2021}, and to the decoupling of row-crop cover from
artificial drainage and of surface from subsurface contributions
\parencite{BolandBrien2014}. Streamflow prediction in tiled landscapes is
itself difficult, raising the unresolved question of whether curve-number or
more process-based methods are required
\parencite{Wessel2022,Thomas2016,Rahman2014}. Whether a representative
catchment scale transfers elsewhere must be established against additional
watersheds \parencite{Basu2010,Thomas2016}, and newly proposed physically based
tools such as the Advanced Terrestrial Simulator and HYDRUS-2D warrant scaling
and transferability tests across soils, topography, land-use change, and
climate \parencite{Rathore2025,Naderi2025,Basu2010}.

\subsection{Channel-network, connectivity, and drain-to-stream processes}

Interactions between drains and the receiving channel network are a frontier in
their own right. Open questions include the fate and transport of tile effluent
after it reaches surface drains \parencite{Clement2017}, drain--river and
drain--ditch interactions \parencite{Grangeon2021}, and field-scale monitoring
of how tile flow affects downstream discharge, stream power, and channel
geomorphology \parencite{Zhou2025,Zhou2013}.

Channel-network, in-stream, and sediment processes are underdeveloped in
watershed models. Networks should be modelled as interdependent reaches with
explicit flow paths and travel times \parencite{Giraud1997,Williamson2025};
in-stream biogeochemistry and hillslope--riparian interaction need stronger
representation \parencite{Schilling2012b,Basu2010}; and gully, streambank, and
non-field sediment sources require mechanistic treatment within catchment-scale
sediment budgets \parencite{Maalim2013,Dalzell2018,Walling2002}.

Connectivity metrics are increasingly flagged as management-relevant. Mapping
surface ditches may be more efficient than locating in-field drains, and future
phosphorus-risk models should incorporate ditch connectivity and lag times
\parencite{Moloney2020,Opoku2024}, supported by improved connectivity
algorithms and monitoring-based thresholds for lowland catchments
\parencite{Calsamiglia2018}. Finally, drained closed depressions and potholes
are singled out as poorly understood: they demand process study at additional
locations, quantification of runoff--tile--antecedent-moisture relationships,
and data on blind-inlet longevity
\parencite{Martin2019,Williams2020,Feyereisen2015b}. Across these threads, the
common gap is a lack of coupled, connectivity-aware representations linking
individual drains to the reaches and depressions that ultimately govern
watershed water quality.
